\section*{图目录（原书页码）}\phantomsection\label{front:figures}
\addcontentsline{toc}{section}{图目录（原书页码）}
\begingroup\small
\begin{longtable}{@{}p{.86\textwidth}r@{}}
\hyperref[fig:1.2.1]{1.2.1 例1.2.7的飞镖靶} & \hyperref[fig:1.2.1]{7} \\
\hyperref[fig:1.2.2]{1.2.2 平均值的直方图} & \hyperref[fig:1.2.2]{16} \\
\hyperref[fig:1.5.1]{1.5.1 例1.5.2的累积分布函数} & \hyperref[fig:1.5.1]{25} \\
\hyperref[fig:1.5.2]{1.5.2 几何分布的累积分布函数，$p=.3$} & \hyperref[fig:1.5.2]{27} \\
\hyperref[fig:1.6.1]{1.6.1 逻辑斯蒂曲线下的面积} & \hyperref[fig:1.6.1]{30} \\
\hyperref[fig:2.1.1]{2.1.1 例2.1.2的变换} & \hyperref[fig:2.1.1]{40} \\
\hyperref[fig:2.1.2]{2.1.2 递增与非递增的累积分布函数}\footnote{译注：原图目录写“increasing and nonincreasing cdfs”，这里照译；累积分布函数应为非递减函数.} & \hyperref[fig:2.1.2]{44} \\
\hyperref[fig:2.3.1]{2.3.1 指数密度} & \hyperref[fig:2.3.1]{49} \\
\hyperref[fig:2.3.2]{2.3.2 具有相同矩的两个概率密度函数} & \hyperref[fig:2.3.2]{53} \\
\hyperref[fig:2.3.3]{2.3.3 二项分布的Poisson近似} & \hyperref[fig:2.3.3]{56} \\
\hyperref[fig:3.3.1]{3.3.1 标准正态密度} & \hyperref[fig:3.3.1]{86} \\
\hyperref[fig:3.3.2]{3.3.2 二项分布的正态近似} & \hyperref[fig:3.3.2]{86} \\
\hyperref[fig:3.3.3]{3.3.3 Beta密度} & \hyperref[fig:3.3.3]{88} \\
\hyperref[fig:3.3.4]{3.3.4 对称Beta密度} & \hyperref[fig:3.3.4]{88} \\
\hyperref[fig:3.3.5]{3.3.5 标准正态密度与Cauchy密度} & \hyperref[fig:3.3.5]{89} \\
\hyperref[fig:3.3.6]{3.3.6 对数正态与Gamma概率密度函数} & \hyperref[fig:3.3.6]{91} \\
\hyperref[fig:3.5.1]{3.5.1 位置密度} & \hyperref[fig:3.5.1]{96} \\
\hyperref[fig:3.5.2]{3.5.2 指数位置密度} & \hyperref[fig:3.5.2]{98} \\
\hyperref[fig:3.5.3]{3.5.3 同一尺度族的成员} & \hyperref[fig:3.5.3]{98} \\
\hyperref[fig:3.5.4]{3.5.4 位置–尺度族} & \hyperref[fig:3.5.4]{99} \\
\hyperref[fig:4.1.1]{4.1.1 积分区域（一）} & \hyperref[fig:4.1.1]{120} \\
\hyperref[fig:4.5.1]{4.5.1 积分区域（二）} & \hyperref[fig:4.5.1]{142} \\
\hyperref[fig:4.5.2]{4.5.2 积分区域（三）} & \hyperref[fig:4.5.2]{143} \\
\hyperref[fig:4.7.1]{4.7.1 凸函数} & \hyperref[fig:4.7.1]{155} \\
\hyperref[fig:4.7.2]{4.7.2 Jensen不等式} & \hyperref[fig:4.7.2]{156} \\
\hyperref[fig:5.4.1]{5.4.1 例5.4.7中$f_{R,V}(r,v)>0$的区域} & \hyperref[fig:5.4.1]{189} \\
\hyperref[fig:5.6.1]{5.6.1 指数概率密度函数的直方图} & \hyperref[fig:5.6.1]{202} \\
\hyperref[fig:5.6.2]{5.6.2 Poisson样本方差的直方图} & \hyperref[fig:5.6.2]{204} \\
\hyperref[fig:5.6.3]{5.6.3 接受–拒绝抽样中的Beta分布} & \hyperref[fig:5.6.3]{205} \\
\hyperref[fig:7.3.1]{7.3.1 二项分布的均方误差比较} & \hyperref[fig:7.3.1]{269} \\
\hyperref[fig:7.3.2]{7.3.2 方差估计量的风险函数} & \hyperref[fig:7.3.2]{284} \\
\hyperref[fig:8.2.1]{8.2.1 似然比检验统计量} & \hyperref[fig:8.2.1]{305} \\
\hyperref[fig:8.3.1]{8.3.1 例8.3.2的功效函数} & \hyperref[fig:8.3.1]{311} \\
\hyperref[fig:8.3.2]{8.3.2 例8.3.3的功效函数} & \hyperref[fig:8.3.2]{312} \\
\hyperref[fig:8.3.3]{8.3.3 例8.3.19中三个检验的功效函数} & \hyperref[fig:8.3.3]{319} \\
\hyperref[fig:8.3.4]{8.3.4 例8.3.31中检验的风险函数} & \hyperref[fig:8.3.4]{325} \\
\hyperref[fig:9.2.1]{9.2.1 置信区间与接受域的关系} & \hyperref[fig:9.2.1]{340} \\
\hyperref[fig:9.2.2]{9.2.2 例9.2.3的接受域与置信区间} & \hyperref[fig:9.2.2]{342} \\
\hyperref[fig:9.2.3]{9.2.3 例9.2.16中的可信区间与置信区间} & \hyperref[fig:9.2.3]{353} \\
\hyperref[fig:9.2.4]{9.2.4 例9.2.16中的可信概率与区间} & \hyperref[fig:9.2.4]{354} \\
\hyperref[fig:9.2.5]{9.2.5 例9.2.16中可信区间的覆盖概率} & \hyperref[fig:9.2.5]{355} \\
\hyperref[fig:9.3.1]{9.3.1 例9.2.16中的三个区间估计量} & \hyperref[fig:9.3.1]{363} \\
\hyperref[fig:10.1.1]{10.1.1 Gamma均值估计量的渐近相对效率} & \hyperref[fig:10.1.1]{386} \\
\hyperref[fig:10.3.1]{10.3.1 Poisson似然的似然比检验} & \hyperref[fig:10.3.1]{396} \\
\hyperref[fig:10.4.1]{10.4.1 二项比例的似然比检验区间} & \hyperref[fig:10.4.1]{406} \\
\hyperref[fig:10.4.2]{10.4.2 名义置信水平为$.9$的二项置信程序的覆盖概率} & \hyperref[fig:10.4.2]{407} \\
\hyperref[fig:11.3.1]{11.3.1 残差平方和所度量的竖直距离} & \hyperref[fig:11.3.1]{438} \\
\hyperref[fig:11.3.2]{11.3.2 最优线性无偏估计量的几何描述} & \hyperref[fig:11.3.2]{443} \\
\hyperref[fig:11.3.3]{11.3.3 Scheffé带、$t$区间与Bonferroni区间} & \hyperref[fig:11.3.3]{455} \\
\hyperref[fig:12.2.1]{12.2.1 正交最小二乘所最小化的距离} & \hyperref[fig:12.2.1]{470} \\
\hyperref[fig:12.2.2]{12.2.2 三条回归直线} & \hyperref[fig:12.2.2]{472} \\
\hyperref[fig:12.2.3]{12.2.3 Creasy–Williams $F$统计量} & \hyperref[fig:12.2.3]{478} \\
\hyperref[fig:12.3.1]{12.3.1 挑战者号数据的逻辑斯蒂曲线} & \hyperref[fig:12.3.1]{481} \\
\hyperref[fig:12.4.1]{12.4.1 最小二乘、最小绝对偏差与M估计拟合} & \hyperref[fig:12.4.1]{485} \\
\end{longtable}\endgroup

% 用户授权疑误脚注：2026-11-04
